\documentclass[11pt]{amsart}
\usepackage{geometry}                % See geometry.pdf to learn the layout options. There are lots.
\geometry{letterpaper}                   % ... or a4paper or a5paper or ... 
%\geometry{landscape}                % Activate for for rotated page geometry
%\usepackage[parfill]{parskip}    % Activate to begin paragraphs with an empty line rather than an indent
\usepackage{graphicx}
\usepackage{amssymb}
\usepackage{epstopdf}
\DeclareGraphicsRule{.tif}{png}{.png}{`convert #1 `dirname #1`/`basename #1 .tif`.png}

\title{Notes on accurate bottom BC on the buoyancy equation}
\author{GPT 5.5 with  prompting from TWNH }
\date{26Jun26}                                           % Activate to display a given date or no date

\begin{document}
\maketitle
%\section{}
%\subsection{}

\section{Boundary-reduced solution using a Dirichlet-to-Neumann map}

Consider the steady buoyancy problem in the domain
\[
\Omega=\{(x,y,z):-H(x,y)\le z\le 0\}.
\]
For simplicity, take the interior equation to contain only vertical diffusion,
\[
\varepsilon^2 \partial_z\left(\kappa \partial_z b\right)
+\gamma(B-b)=0,
\qquad -H(x,y)<z<0,
\]
with inhomogeneous surface Dirichlet condition
\[
b(x,y,0)=b_s(x,y),
\]
and no-flux condition on the sloping bottom
\[
\left[
b_z+\delta^2 \nabla_h b\cdot \nabla_h H
\right]_{z=-H(x,y)}=0.
\]
Here
\[
\nabla_h=(\partial_x,\partial_y).
\]

The difficulty is that the bottom boundary condition couples neighboring water columns through the term
\[
\nabla_h b\cdot \nabla_h H.
\]
A convenient way to isolate this coupling is to reduce the problem to an equation posed only on the bottom boundary.

Define the bottom trace
\[
\beta(x,y)=b(x,y,-H(x,y)).
\]

The horizontal gradient of the bottom trace is not the same as the horizontal gradient of \(b\) evaluated at fixed \(z\). By the chain rule,
\[
\nabla_h \beta
=
\nabla_h b|_{z=-H}
+
b_z|_{z=-H}\nabla_h(-H).
\]
Since
\[
\nabla_h(-H)=-\nabla_h H,
\]
we have
\[
\nabla_h \beta
=
\nabla_h b|_{z=-H}
-
b_z|_{z=-H}\nabla_h H.
\]
Therefore,
\[
\nabla_h b|_{z=-H}
=
\nabla_h \beta
+
b_z|_{z=-H}\nabla_h H.
\]

Substituting this into the bottom no-flux condition gives
\[
b_z+\delta^2
\left(
\nabla_h \beta
+
b_z\nabla_h H
\right)\cdot \nabla_h H=0.
\]
Therefore,
\[
b_z
+\delta^2 \nabla_h \beta\cdot \nabla_h H
+\delta^2 b_z |\nabla_h H|^2=0.
\]
Equivalently,
\[
\boxed{
\left(1+\delta^2 |\nabla_h H|^2\right)b_z|_{z=-H}
+
\delta^2 \nabla_h H\cdot \nabla_h \beta
=0.
}
\]

Now solve the vertical column problem at each horizontal location with prescribed surface value \(b_s(x,y)\) and prescribed bottom value \(\beta(x,y)\):
\[
\varepsilon^2 \partial_z\left(\kappa \partial_z b\right)
+\gamma(B-b)=0,
\qquad -H<z<0,
\]
\[
b(0)=b_s,
\qquad
b(-H)=\beta.
\]
This vertical solve defines a Dirichlet-to-Neumann map. Namely, it gives the bottom vertical derivative as
\[
b_z(-H)=\Lambda[\beta]+r.
\]
Here \(\Lambda\) is the columnwise Dirichlet-to-Neumann operator mapping the imposed bottom trace \(\beta\) to the resulting bottom normal derivative, while \(r\) is a known contribution coming from the surface value \(b_s\) and the interior forcing \(B\).

Substituting
\[
b_z(-H)=\Lambda[\beta]+r
\]
into the transformed bottom condition gives the reduced horizontal equation for the unknown bottom trace:
\[
\boxed{
\left(1+\delta^2 |\nabla_h H|^2\right)
\left(\Lambda[\beta]+r\right)
+
\delta^2 \nabla_h H\cdot \nabla_h \beta
=0.
}
\]
Once this equation has been solved for \(\beta(x,y)\), the full buoyancy field is reconstructed by solving the vertical column problem with
\[
b(0)=b_s(x,y),
\qquad
b(-H)=\beta(x,y).
\]

\subsection{Constant diffusivity example}

Assume
\[
\kappa=\kappa_0
\]
and define
\[
\psi=\frac{1}{\varepsilon}\sqrt{\frac{\gamma}{\kappa_0}}.
\]
Then the vertical equation becomes
\[
b_{zz}-\psi^2 b=-\psi^2 B.
\]

For clarity, first consider the homogeneous case \(B=0\). Then
\[
b_{zz}-\psi^2 b=0,
\]
with
\[
b(0)=b_s,
\qquad
b(-H)=\beta.
\]
The solution is
\[
b(z)
=
\beta
\frac{\sinh(-\psi z)}{\sinh(\psi H)}
+
b_s
\frac{\sinh(\psi(z+H))}{\sinh(\psi H)}.
\]
Differentiating,
\[
b_z(z)
=
-\psi\beta
\frac{\cosh(-\psi z)}{\sinh(\psi H)}
+
\psi b_s
\frac{\cosh(\psi(z+H))}{\sinh(\psi H)}.
\]
Evaluating at the bottom \(z=-H\) gives
\[
b_z(-H)
=
-\psi \coth(\psi H)\beta
+
\psi \operatorname{csch}(\psi H)b_s.
\]
Therefore, in this case,
\[
\Lambda[\beta]
=
-\psi \coth(\psi H)\beta,
\]
and
\[
r=
\psi \operatorname{csch}(\psi H)b_s.
\]

The reduced bottom-trace equation is therefore
\[
\left(1+\delta^2 |\nabla_h H|^2\right)
\left[
-\psi \coth(\psi H)\beta
+
\psi \operatorname{csch}(\psi H)b_s
\right]
+
\delta^2 \nabla_h H\cdot \nabla_h \beta
=0.
\]
Equivalently,
\[
\boxed{
\delta^2 \nabla_h H\cdot \nabla_h \beta
-
\left(1+\delta^2 |\nabla_h H|^2\right)
\psi \coth(\psi H)\beta
=
-
\left(1+\delta^2 |\nabla_h H|^2\right)
\psi \operatorname{csch}(\psi H)b_s.
}
\]

After solving this horizontal equation for \(\beta(x,y)\), the full solution is recovered from
\[
\boxed{
b(z)
=
\beta
\frac{\sinh(-\psi z)}{\sinh(\psi H)}
+
b_s
\frac{\sinh(\psi(z+H))}{\sinh(\psi H)}.
}
\]

\subsection{Including an interior buoyancy profile \(B\)}

When \(B\neq 0\), the vertical equation is
\[
b_{zz}-\psi^2 b=-\psi^2 B(z;x,y).
\]
The solution may be written as
\[
b(z)
=
b_{\mathrm{hom}}(z)
+
b_{\mathrm{part}}(z),
\]
where
\[
b_{\mathrm{hom}}(z)
=
\beta
\frac{\sinh(-\psi z)}{\sinh(\psi H)}
+
b_s
\frac{\sinh(\psi(z+H))}{\sinh(\psi H)}
\]
and \(b_{\mathrm{part}}\) solves
\[
b_{\mathrm{part},zz}-\psi^2 b_{\mathrm{part}}
=
-\psi^2 B,
\]
with homogeneous endpoint conditions
\[
b_{\mathrm{part}}(0)=0,
\qquad
b_{\mathrm{part}}(-H)=0.
\]
Then
\[
b_z(-H)
=
-\psi \coth(\psi H)\beta
+
\psi \operatorname{csch}(\psi H)b_s
+
b_{\mathrm{part},z}(-H).
\]
Thus,
\[
\Lambda[\beta]
=
-\psi \coth(\psi H)\beta,
\]
and
\[
r=
\psi \operatorname{csch}(\psi H)b_s
+
b_{\mathrm{part},z}(-H).
\]

Using the Dirichlet Green's function \(G_D(z,\xi)\) on \([-H,0]\), satisfying
\[
\left(\partial_z^2-\psi^2\right)G_D(z,\xi)
=
\delta(z-\xi),
\]
\[
G_D(0,\xi)=0,
\qquad
G_D(-H,\xi)=0,
\]
we may write
\[
b_{\mathrm{part}}(z)
=
-\psi^2 \int_{-H}^{0} G_D(z,\xi)B(\xi)\,d\xi.
\]
Therefore,
\[
b_{\mathrm{part},z}(-H)
=
-\psi^2 \int_{-H}^{0} \partial_z G_D(-H,\xi)B(\xi)\,d\xi.
\]

For constant \(\psi\), the Dirichlet Green's function is
\[
G_D(z,\xi)
=
-\frac{1}{\psi\sinh(\psi H)}
\begin{cases}
\sinh\left(\psi(z+H)\right)\sinh(-\psi \xi),
& z\le \xi,\\[4pt]
\sinh\left(\psi(\xi+H)\right)\sinh(-\psi z),
& z\ge \xi.
\end{cases}
\]
At the bottom, \(z=-H\le \xi\), so
\[
\partial_z G_D(-H,\xi)
=
-\frac{\sinh(-\psi \xi)}{\sinh(\psi H)}.
\]
Hence
\[
b_{\mathrm{part},z}(-H)
=
\psi^2
\int_{-H}^{0}
\frac{\sinh(-\psi \xi)}{\sinh(\psi H)}
B(\xi)\,d\xi.
\]
Therefore,
\[
\boxed{
r
=
\psi \operatorname{csch}(\psi H)b_s
+
\psi^2
\int_{-H}^{0}
\frac{\sinh(-\psi \xi)}{\sinh(\psi H)}
B(\xi)\,d\xi.
}
\]

The reduced equation with interior forcing is
\[
\boxed{
\left(1+\delta^2 |\nabla_h H|^2\right)
\left[
-\psi \coth(\psi H)\beta
+
r
\right]
+
\delta^2 \nabla_h H\cdot \nabla_h \beta
=0.
}
\]

Equivalently,
\[
\boxed{
\delta^2 \nabla_h H\cdot \nabla_h \beta
-
\left(1+\delta^2 |\nabla_h H|^2\right)
\psi \coth(\psi H)\beta
=
-
\left(1+\delta^2 |\nabla_h H|^2\right)r.
}
\]

\subsection{Interpretation}

The reduction has converted the original sloping-bottom boundary-value problem into two pieces:

\begin{enumerate}
    \item A horizontal equation for the bottom trace \(\beta(x,y)\):
    \[
    \left(1+\delta^2 |\nabla_h H|^2\right)
    \left(\Lambda[\beta]+r\right)
    +
    \delta^2 \nabla_h H\cdot \nabla_h \beta
    =0.
    \]

    \item A vertical reconstruction problem:
    \[
    \varepsilon^2 \partial_z(\kappa b_z)+\gamma(B-b)=0,
    \qquad
    b(0)=b_s,
    \qquad
    b(-H)=\beta.
    \]
\end{enumerate}

The only horizontal coupling appears in the bottom-trace equation. In the limit \(\delta\to 0\), the equation reduces to
\[
\Lambda[\beta]+r=0,
\]
which enforces
\[
b_z(-H)=0
\]
column by column.

For the **Dirichlet-to-Neumann reduction / boundary trace approach**, the standard references are from elliptic PDE theory and domain-decomposition theory. The idea is usually called the **Steklov–Poincaré operator**, **Poincaré–Steklov operator**, or **Dirichlet-to-Neumann map**.

\section{Solving the reduced bottom-trace equation for $\beta$}

We consider a closed horizontal domain \(D\subset\mathbb{R}^2\), with no lateral
inflow boundary prescribed for the reduced bottom-trace problem. The surface
boundary condition at \(z=0\) specifies the buoyancy:
\[
b(x,y,0)=b_s(x,y).
\]
The bottom trace is
\[
\beta(x,y)=b(x,y,-H(x,y)).
\]

After applying the Dirichlet-to-Neumann reduction, the sloping-bottom no-flux
condition gives a first-order horizontal equation for \(\beta\). In the constant
diffusivity example this equation has the form
\[
\delta^2 \nabla_h H\cdot \nabla_h \beta
-
\left(1+\delta^2 |\nabla_h H|^2\right)
\psi \coth(\psi H)\beta
=
-
\left(1+\delta^2 |\nabla_h H|^2\right)r,
\]
where
\[
\psi=\frac{1}{\varepsilon}\sqrt{\frac{\gamma}{\kappa_0}},
\]
and \(r(x,y)\) is a known forcing contribution determined by the prescribed
surface value \(b_s\) and the interior buoyancy profile \(B\).

Equivalently, define
\[
\mathbf{a}(x,y)=-\delta^2 \nabla_h H(x,y),
\]
\[
\lambda(x,y)
=
\left(1+\delta^2 |\nabla_h H|^2\right)\psi\coth(\psi H),
\]
and
\[
q(x,y)
=
\left(1+\delta^2 |\nabla_h H|^2\right)r(x,y).
\]
Then the reduced bottom-trace equation is
\[
\boxed{
\mathbf{a}\cdot\nabla_h\beta+\lambda\beta=q.
}
\]
Since
\[
\lambda(x,y)>0,
\]
this is a linear advection--reaction equation with positive reaction.

\subsection{Characteristics}

The characteristic curves \(\mathbf{X}(s)\) satisfy
\[
\frac{d\mathbf{X}}{ds}=\mathbf{a}(\mathbf{X})
=
-\delta^2 \nabla_h H(\mathbf{X}).
\]
Therefore the characteristics are downhill gradient-flow curves of the
topography \(H\). Indeed,
\[
\frac{d}{ds}H(\mathbf{X}(s))
=
\nabla_h H\cdot \frac{d\mathbf{X}}{ds}
=
-\delta^2 |\nabla_h H|^2
\le 0.
\]
Thus \(H\) is non-increasing along characteristics. Equality holds only at
critical points of \(H\), where
\[
\nabla_h H=0.
\]

Along a characteristic,
\[
\frac{d}{ds}\beta(\mathbf{X}(s))
=
\mathbf{a}(\mathbf{X}(s))\cdot\nabla_h\beta,
\]
and the PDE becomes the scalar ODE
\[
\boxed{
\frac{d\beta}{ds}+\lambda(\mathbf{X}(s))\beta
=
q(\mathbf{X}(s)).
}
\]
The solution along a characteristic from \(s_0\) to \(s\) is
\[
\boxed{
\beta(s)
=
\beta(s_0)
\exp\left[
-\int_{s_0}^{s}\lambda(\mathbf{X}(\tau))\,d\tau
\right]
+
\int_{s_0}^{s}
\exp\left[
-\int_{\sigma}^{s}\lambda(\mathbf{X}(\tau))\,d\tau
\right]
q(\mathbf{X}(\sigma))\,d\sigma.
}
\]

\subsection{Closed domains and critical points}

In a closed domain, there is no external inflow boundary for the first-order
equation. Instead, characteristics originate from local maxima of \(H\), or from
critical sets where
\[
\nabla_h H=0.
\]
At such points the advective term vanishes and the reduced equation becomes
purely algebraic:
\[
\lambda\beta=q.
\]
Therefore, at a critical point,
\[
\boxed{
\beta=\frac{q}{\lambda}.
}
\]
For the present coefficients,
\[
q=
\left(1+\delta^2|\nabla_h H|^2\right)r,
\]
and
\[
\lambda=
\left(1+\delta^2|\nabla_h H|^2\right)\psi\coth(\psi H).
\]
Thus, at any critical point of \(H\), where \(|\nabla_h H|=0\),
\[
\boxed{
\beta
=
\frac{r}{\psi\coth(\psi H)}
=
\frac{r}{\psi}\tanh(\psi H).
}
\]

This value supplies the upstream value for characteristics that emanate from
local maxima of \(H\). The value of \(\beta\) is then propagated downhill along
the gradient-flow curves using the characteristic formula.

\subsection{Topology of the characteristic field}

Because the characteristic velocity is
\[
\mathbf{a}=-\delta^2\nabla_h H,
\]
the characteristic curves are gradient-flow curves. If \(H\) is smooth and
\(\nabla_h H\) is finite, then solutions of
\[
\frac{d\mathbf{X}}{ds}=-\delta^2\nabla_h H
\]
are unique away from critical points. Therefore, two distinct characteristics
cannot cross at a regular point.

However, the characteristic topology may still be complicated:

\begin{itemize}
    \item characteristics descend from local maxima of \(H\);
    \item they can approach saddles or local minima;
    \item separatrices divide regions fed by different maxima;
    \item different characteristic families may meet at critical points;
    \item if upstream values differ across separatrices, \(\beta\) may develop
    sharp jumps or weak discontinuities there.
\end{itemize}

The positive reaction term \(\lambda\beta\) damps sensitivity to upstream data.
Differences in upstream values decay approximately like
\[
\exp\left[-\int \lambda\,ds\right]
\]
along characteristics. However, without horizontal diffusion, the equation does
not smooth across characteristics or across separatrices.

\subsection{Numerical solution by upwind discretization}

For computation, it is often preferable to solve
\[
\mathbf{a}\cdot\nabla_h\beta+\lambda\beta=q
\]
directly using an upwind finite-difference or finite-volume method rather than
explicitly tracing all characteristics.

On a Cartesian grid, write
\[
\mathbf{a}=(a_x,a_y).
\]
Then
\[
a_x\partial_x\beta+a_y\partial_y\beta+\lambda\beta=q.
\]
Use upwind differences:
\[
\partial_x\beta\approx
\begin{cases}
\dfrac{\beta_{i,j}-\beta_{i-1,j}}{\Delta x}, & a_x>0,\\[10pt]
\dfrac{\beta_{i+1,j}-\beta_{i,j}}{\Delta x}, & a_x<0,
\end{cases}
\]
and
\[
\partial_y\beta\approx
\begin{cases}
\dfrac{\beta_{i,j}-\beta_{i,j-1}}{\Delta y}, & a_y>0,\\[10pt]
\dfrac{\beta_{i,j+1}-\beta_{i,j}}{\Delta y}, & a_y<0.
\end{cases}
\]
This produces a sparse linear system
\[
A\boldsymbol{\beta}=\mathbf{q}.
\]
The reaction coefficient \(\lambda>0\) makes the problem well-conditioned
relative to a pure advection problem.

In a closed domain, one does not impose an external inflow condition. Instead,
the behavior near critical points of \(H\) determines the upstream values
through the algebraic balance
\[
\beta=\frac{q}{\lambda}.
\]
A robust discretization should therefore treat cells where
\[
|\nabla_h H|\approx 0
\]
as reaction-dominated cells and set, or strongly relax toward,
\[
\beta\approx \frac{q}{\lambda}.
\]

\subsection{Small-$\delta$ expansion}

Since the horizontal advection term is proportional to \(\delta^2\), it may be
useful to solve the reduced equation perturbatively. The equation is
\[
-\delta^2\nabla_h H\cdot\nabla_h\beta+\lambda\beta=q,
\]
where
\[
\lambda=
\left(1+\delta^2|\nabla_h H|^2\right)\psi\coth(\psi H),
\]
and
\[
q=
\left(1+\delta^2|\nabla_h H|^2\right)r.
\]

Expand
\[
\beta=\beta^{(0)}+\delta^2\beta^{(1)}+\cdots,
\]
\[
\lambda=\lambda_0+\delta^2\lambda_1+\cdots,
\]
\[
q=q_0+\delta^2 q_1+\cdots.
\]
At leading order,
\[
\lambda_0\beta^{(0)}=q_0,
\]
where
\[
\lambda_0=\psi\coth(\psi H),
\qquad
q_0=r.
\]
Thus
\[
\boxed{
\beta^{(0)}
=
\frac{r}{\psi\coth(\psi H)}
=
\frac{r}{\psi}\tanh(\psi H).
}
\]

At order \(\delta^2\),
\[
-\nabla_h H\cdot\nabla_h\beta^{(0)}
+\lambda_0\beta^{(1)}
+\lambda_1\beta^{(0)}
=
q_1.
\]
Therefore,
\[
\boxed{
\beta^{(1)}
=
\frac{
q_1
+
\nabla_h H\cdot\nabla_h\beta^{(0)}
-
\lambda_1\beta^{(0)}
}{
\lambda_0
}.
}
\]
This perturbative approach avoids characteristic-tracing issues and is natural
when \(\delta^2\ll 1\).

\subsection{Regularization by horizontal diffusion}

The first-order equation
\[
\mathbf{a}\cdot\nabla_h\beta+\lambda\beta=q
\]
does not smooth across characteristics. If the characteristic topology contains
saddles, separatrices, or reconnecting basins, the solution may develop sharp
gradients or weak discontinuities.

A more robust model is obtained by retaining a small horizontal diffusion term
in the original buoyancy equation. The reduced bottom-trace equation then becomes
an advection--diffusion--reaction problem, schematically
\[
-\nabla_h\cdot(D_{\mathrm{eff}}\nabla_h\beta)
+
\mathbf{a}\cdot\nabla_h\beta
+
\lambda\beta
=
q.
\]
This is elliptic or advection--diffusion--reaction type and no longer requires
delicate characteristic boundary conditions. It also smooths across separatrices.

\subsection{Summary}

The bottom-trace equation in a closed domain is
\[
\boxed{
\mathbf{a}\cdot\nabla_h\beta+\lambda\beta=q,
\qquad
\mathbf{a}=-\delta^2\nabla_h H,
}
\]
with
\[
\lambda=
\left(1+\delta^2|\nabla_h H|^2\right)\psi\coth(\psi H),
\qquad
q=
\left(1+\delta^2|\nabla_h H|^2\right)r.
\]
It is a linear advection--reaction equation. The characteristics are downhill
gradient-flow curves of \(H\):
\[
\frac{d\mathbf{X}}{ds}=-\delta^2\nabla_h H.
\]
At critical points of \(H\), the equation reduces to
\[
\beta=\frac{q}{\lambda}.
\]
In practice, the equation can be solved either by characteristics, by an upwind
finite-volume/finite-difference discretization, or perturbatively in
\(\delta^2\). If separatrices or reconnecting characteristic basins cause
singular behavior, horizontal diffusion should be retained, yielding an
advection--diffusion--reaction equation for \(\beta\).

\end{document}  